\documentclass[10pt,a4paper]{amsart}
\usepackage{amsmath,amssymb,amsthm,graphicx}
% Retained minimal preamble for the Stacks Project chapter sites-cohomology.tex (Cohomology on Sites).
% Verbatim excerpt of the chapter preamble: Xy-pic setup (native preamble lines 14-18) and the native
% theorem declarations with the mathematical macros (native preamble lines 154-217). Copied mechanically;
% nothing added, removed or redefined. The wrapper class and the amsmath/amssymb/amsthm/graphicx base
% packages are supplied by the assembler, so the class/externaldocument part of the native preamble is omitted.
% For commutative diagrams we use Xy-pic
\usepackage[all]{xy}

% We use 2cell for 2-commutative diagrams.
\xyoption{2cell}
% Theorem environments.
%
\theoremstyle{plain}
\newtheorem{theorem}[subsection]{Theorem}
\newtheorem{proposition}[subsection]{Proposition}
\newtheorem{lemma}[subsection]{Lemma}

\theoremstyle{definition}
\newtheorem{definition}[subsection]{Definition}
\newtheorem{example}[subsection]{Example}
\newtheorem{exercise}[subsection]{Exercise}
\newtheorem{situation}[subsection]{Situation}

\theoremstyle{remark}
\newtheorem{remark}[subsection]{Remark}
\newtheorem{remarks}[subsection]{Remarks}

\numberwithin{equation}{subsection}

% Macros
%
\def\lim{\mathop{\mathrm{lim}}\nolimits}
\def\colim{\mathop{\mathrm{colim}}\nolimits}
\def\Spec{\mathop{\mathrm{Spec}}}
\def\Hom{\mathop{\mathrm{Hom}}\nolimits}
\def\Ext{\mathop{\mathrm{Ext}}\nolimits}
\def\SheafHom{\mathop{\mathcal{H}\!\mathit{om}}\nolimits}
\def\SheafExt{\mathop{\mathcal{E}\!\mathit{xt}}\nolimits}
\def\Sch{\mathit{Sch}}
\def\Mor{\mathop{\mathrm{Mor}}\nolimits}
\def\Ob{\mathop{\mathrm{Ob}}\nolimits}
\def\Sh{\mathop{\mathit{Sh}}\nolimits}
\def\NL{\mathop{N\!L}\nolimits}
\def\CH{\mathop{\mathrm{CH}}\nolimits}
\def\proetale{{pro\text{-}\acute{e}tale}}
\def\etale{{\acute{e}tale}}
\def\QCoh{\mathit{QCoh}}
\def\Ker{\mathop{\mathrm{Ker}}}
\def\Im{\mathop{\mathrm{Im}}}
\def\Coker{\mathop{\mathrm{Coker}}}
\def\Coim{\mathop{\mathrm{Coim}}}

% Boxtimes
%
\DeclareMathSymbol{\boxtimes}{\mathbin}{AMSa}{"02}

%
% Macros for moduli stacks/spaces
%
\def\QCohstack{\mathcal{QC}\!\mathit{oh}}
\def\Cohstack{\mathcal{C}\!\mathit{oh}}
\def\Spacesstack{\mathcal{S}\!\mathit{paces}}
\def\Quotfunctor{\mathrm{Quot}}
\def\Hilbfunctor{\mathrm{Hilb}}
\def\Curvesstack{\mathcal{C}\!\mathit{urves}}
\def\Polarizedstack{\mathcal{P}\!\mathit{olarized}}
\def\Complexesstack{\mathcal{C}\!\mathit{omplexes}}
% \Pic is the operator that assigns to X its picard group, usage \Pic(X)
% \Picardstack_{X/B} denotes the Picard stack of X over B
% \Picardfunctor_{X/B} denotes the Picard functor of X over B
\def\Pic{\mathop{\mathrm{Pic}}\nolimits}
\def\Picardstack{\mathcal{P}\!\mathit{ic}}
\def\Picardfunctor{\mathrm{Pic}}
\def\Deformationcategory{\mathcal{D}\!\mathit{ef}}
\usepackage{cleveref}
\crefname{originallemma}{original lemma}{original lemmas}
\Crefname{originallemma}{Original lemma}{Original lemmas}
\makeatletter
\newlabel{sites-modules-section-localize}{{03DH}{source}{}{sourcepointer-1}{}}
\newlabel{sites-modules-section-localize@cref}{{[originallemma][0][]03DH}{[1][0][]source}{}{}{}}
\makeatother
\crefname{originallemma}{original lemma}{original lemmas}
\Crefname{originallemma}{Original lemma}{Original lemmas}
\makeatletter
\newlabel{categories-lemma-yoneda}{{001P}{source}{}{sourcepointer-2}{}}
\newlabel{categories-lemma-yoneda@cref}{{[originallemma][0][]001P}{[1][0][]source}{}{}{}}
\makeatother
\crefname{originallemma}{original lemma}{original lemmas}
\Crefname{originallemma}{Original lemma}{Original lemmas}
\makeatletter
\newlabel{homology-lemma-characterize-injectives}{{0136}{source}{}{sourcepointer-3}{}}
\newlabel{homology-lemma-characterize-injectives@cref}{{[originallemma][0][]0136}{[1][0][]source}{}{}{}}
\makeatother
\crefname{originallemma}{original lemma}{original lemmas}
\Crefname{originallemma}{Original lemma}{Original lemmas}
\makeatletter
\newlabel{sites-modules-lemma-extension-by-zero}{{03DI}{source}{}{sourcepointer-4}{}}
\newlabel{sites-modules-lemma-extension-by-zero@cref}{{[originallemma][0][]03DI}{[1][0][]source}{}{}{}}
\makeatother
\title{Restriction along a monomorphism is surjective on injective modules}
\author{The Stacks Project Authors}
\date{stacks-project:sites-cohomology.tex:local-snapshot-20261011; GNU FDL 1.2 or later}
\begin{document}
\maketitle
\noindent\textit{Editorial note.} All delivered mathematical text is the author's own, copied from the pinned Stacks Project chapter file sites-cohomology.tex; no statement or proof step was rewritten, reordered or removed.\par\medskip
\noindent\textit{Editorial note.} The proof is the author's complete argument: the localisation set-up and the extension by zero are both inside the counted span.\par\medskip
\noindent\textit{Editorial note.} Every labelled reference occurring inside the counted spans is carried in assembly.external\_references as a non-load-bearing navigation pointer (same-chapter ones at their definition site, cross-chapter ones at the invocation inside the counted span), so no load-bearing mathematics is replaced by a stub. No bibliography citation occurs inside the counted spans.\par\medskip
\noindent\textit{Editorial note.} The checked chapter list for same-name collisions is the set of stacks chapters staged in this repository: chow, cohomology, divisors, etale-cohomology, flat, more-algebra, more-morphisms, morphisms, spaces-more-morphisms, varieties together with this card's chapter sites-cohomology.\par\medskip
\noindent\textit{Editorial note.} The preamble is the chapter's own retained minimal preamble (sources/stacks-sites-cohomology/retained-preamble.tex), a verbatim excerpt of the native preamble: the Xy-pic setup and the native theorem declarations with the mathematical macros.\par\medskip
\noindent\textit{Allowed standard prerequisites (editorial).}\begin{itemize}
\item Localization of a ringed site at an object and restriction of modules. Same chapter, the localized site C/U, the morphism j : C/U -> C and the restriction functor on modules. Supplies the functors j\_! and j\textasciicircum{}\{-1\} through which the restriction map is analysed.
\item Injective module over a ringed site and the adjunction between restriction and its left adjoint. Same chapter, injectivity of a module, the adjunction j\_! -- j\textasciicircum{}\{-1\}, and the monomorphism hypothesis on a : U' -> U which makes j\_!j\textasciicircum{}\{-1\} a subfunctor of the identity. Supplies the extension of a section along the monomorphism used to produce a preimage.
\end{itemize}\medskip
\ifdefined\hypertarget\hypertarget{sourcepointer-1}{}\fi
\noindent\textit{Source reference.} Original lemma 03DH: sources/stacks-sites-cohomology/pinned-original-source.tex, source lines 1972--1972. Navigation pointer invoked at proof line 1972; the counted span contains the pointer only and the cited result is not reproduced inside it, so the pointer is registered as a non-load-bearing navigation stub in the established form for Stacks material. No bibliography citation is involved.\par\medskip
\ifdefined\hypertarget\hypertarget{sourcepointer-2}{}\fi
\noindent\textit{Source reference.} Original lemma 001P: sources/stacks-sites-cohomology/pinned-original-source.tex, source lines 1990--1990. Navigation pointer invoked at proof line 1990; the counted span contains the pointer only and the cited result is not reproduced inside it, so the pointer is registered as a non-load-bearing navigation stub in the established form for Stacks material. No bibliography citation is involved.\par\medskip
\ifdefined\hypertarget\hypertarget{sourcepointer-3}{}\fi
\noindent\textit{Source reference.} Original lemma 0136: sources/stacks-sites-cohomology/pinned-original-source.tex, source lines 1992--1992. Navigation pointer invoked at proof line 1992; the counted span contains the pointer only and the cited result is not reproduced inside it, so the pointer is registered as a non-load-bearing navigation stub in the established form for Stacks material. No bibliography citation is involved.\par\medskip
\ifdefined\hypertarget\hypertarget{sourcepointer-4}{}\fi
\noindent\textit{Source reference.} Original lemma 03DI: sources/stacks-sites-cohomology/pinned-original-source.tex, source lines 2003--2003. Navigation pointer invoked at proof line 2003; the counted span contains the pointer only and the cited result is not reproduced inside it, so the pointer is registered as a non-load-bearing navigation stub in the established form for Stacks material. No bibliography citation is involved.\par\medskip
\begin{lemma}
\label{lemma-restriction-along-monomorphism-surjective}
Let $(\mathcal{C}, \mathcal{O})$ be a ringed site. Let $a : U' \to U$ be a
monomorphism in $\mathcal{C}$. Then for any injective $\mathcal{O}$-module
$\mathcal{I}$ the restriction mapping $\mathcal{I}(U) \to \mathcal{I}(U')$
is surjective.
\end{lemma}

\begin{proof}
Let $j : \mathcal{C}/U \to \mathcal{C}$ and
$j' : \mathcal{C}/U' \to \mathcal{C}$ be the localization morphisms
(Modules on Sites, Section \ref{sites-modules-section-localize}).
Since $j_!$ is a left adjoint to restriction we see that
for any sheaf $\mathcal{F}$ of $\mathcal{O}$-modules
$$
\Hom_\mathcal{O}(j_!\mathcal{O}_U, \mathcal{F})
=
\Hom_{\mathcal{O}_U}(\mathcal{O}_U, \mathcal{F}|_U)
=
\mathcal{F}(U)
$$
Similarly, the sheaf $j'_!\mathcal{O}_{U'}$ represents the
functor $\mathcal{F} \mapsto \mathcal{F}(U')$.
Moreover below we describe a canonical map of $\mathcal{O}$-modules
$$
j'_!\mathcal{O}_{U'} \longrightarrow j_!\mathcal{O}_U
$$
which corresponds to the restriction mapping
$\mathcal{F}(U) \to \mathcal{F}(U')$ via Yoneda's lemma
(Categories, Lemma \ref{categories-lemma-yoneda}).
It suffices to prove the displayed map of modules is injective, see
Homology, Lemma \ref{homology-lemma-characterize-injectives}.

\medskip\noindent
To construct our map it suffices to construct a map between the
presheaves which assign to an object $V$ of $\mathcal{C}$ the
$\mathcal{O}(V)$-module
$$
\bigoplus\nolimits_{\varphi' \in \Mor_\mathcal{C}(V, U')} \mathcal{O}(V)
\quad\text{and}\quad
\bigoplus\nolimits_{\varphi \in \Mor_\mathcal{C}(V, U)} \mathcal{O}(V)
$$
see Modules on Sites, Lemma \ref{sites-modules-lemma-extension-by-zero}.
We take the map which maps the summand corresponding to $\varphi'$
to the summand corresponding to $\varphi = a \circ \varphi'$
by the identity map on $\mathcal{O}(V)$. As $a$ is a monomorphism,
this map is injective. As sheafification is exact, the result
follows.
\end{proof}

\end{document}
